% header.tex — xelatex unicode fallbacks for the Zebra book
%
% Purpose: substitute glyphs that the active fonts (Noto Serif / Noto Sans /
% DejaVu Sans Mono) do not provide, so xelatex stops emitting
% "Missing character" warnings and the PDF prints all source characters
% as readable Latin-script substitutes.
%
% Add a new \newunicodechar mapping if a future build complains about
% an additional codepoint.

\usepackage{newunicodechar}
\usepackage{pifont}      % \ding{51} (heavy check), \ding{55} (heavy ballot X)

% --- emoji used as inline markers --------------------------------------------

% ✅ U+2705  WHITE HEAVY CHECK MARK
\newunicodechar{✅}{{\ding{51}}}

% ❌ U+274C  CROSS MARK
\newunicodechar{❌}{{\ding{55}}}

% 🏋 U+1F3CB  PERSON LIFTING WEIGHTS — used as exercise-section marker
\newunicodechar{🏋}{\textbf{[Exercise]}}

% 💡 U+1F4A1  ELECTRIC LIGHT BULB — used as "Why" / "Tip" note marker
\newunicodechar{💡}{\textbf{Tip:}}

% --- emoji that appear inside source-code examples ---------------------------
% (DejaVu Sans Mono lacks emoji glyphs entirely; ascii substitutes keep the
% example readable without changing the source markdown.)

% 👋 U+1F44B  WAVING HAND
\newunicodechar{👋}{[wave]}

% 🌍 U+1F30D  EARTH GLOBE EUROPE-AFRICA
\newunicodechar{🌍}{[earth]}

% 🎉 U+1F389  PARTY POPPER
\newunicodechar{🎉}{[party]}

% --- arrows ------------------------------------------------------------------

% → U+2192  RIGHTWARDS ARROW
% Math-mode \rightarrow is font-independent; \textrightarrow re-emits U+2192
% in the active text font, which loops back into a missing-glyph warning.
\newunicodechar{→}{\ensuremath{\rightarrow}}

% --- CJK in the Strings-and-Unicode chapter ----------------------------------
% The book's Unicode chapter uses literal CJK ("你好世界" — "Hello World" in
% Chinese) inside a code example. The active monospace font has no CJK
% coverage, so each codepoint is mapped to its Hanyu-pinyin romanization.
% Markdown source is unchanged; only the PDF render substitutes.

\newunicodechar{你}{[ni]}
\newunicodechar{好}{[hao]}
\newunicodechar{世}{[shi]}
\newunicodechar{界}{[jie]}
